\section{Other statistics: what the group changes}\label{sec:other}

\begin{table}[H]
\centering\footnotesize
\begin{tabular}{@{}l l l l@{}}
\toprule
 & solo flights $\Sl$ & group time $\G$ & same flights alone $\A$ \\
\midrule
flights (TA-MSD) & 119\,129 & 33\,435 & 29\,611 \\
$H$ (10--1000~s) & 0.899 [0.898, 0.899] & 0.916 [0.915, 0.917] & 0.914 [0.913, 0.915] \\
$H_\mathrm{loc}(100~\mathrm{s})$ & 0.919 [0.919, 0.919] & 0.927 [0.927, 0.928] & 0.932 [0.931, 0.932] \\
$H_\mathrm{loc}(1000~\mathrm{s})$ & 0.837 [0.837, 0.838] & 0.893 [0.891, 0.894] & 0.860 [0.859, 0.862] \\
time in climb (\%) & 30.8 [30.7, 30.9] & 33.0 [32.8, 33.2] & 29.1 [28.9, 29.2] \\
searches ending in a climb (\%) & 43.3 [43.2, 43.4] & 49.5 [49.2, 49.7] & 45.9 [45.7, 46.1] \\
climb rate (m/s) & 1.22 [1.22, 1.22] & 1.26 [1.25, 1.27] & 1.33 [1.32, 1.34] \\
glide speed (km/h) & 30.1 [30.1, 30.2] & 33.3 [33.1, 33.5] & 32.8 [32.6, 32.9] \\
glide ratio & 10.4 [10.4, 10.4] & 12.9 [12.8, 13.0] & 11.3 [11.2, 11.4] \\
cross-country speed (km/h) & 17.2 [17.2, 17.3] & 20.0 [19.8, 20.3] & 20.0 [19.8, 20.1] \\
\bottomrule
\end{tabular}

\caption{Pooled statistics of S, G and A (paragliders), with 5--95\% bands (blocks of the note). TA-MSD on loggers of at most 10~s, phases on 1~Hz flights.}
\label{tab:raw}
\end{table}

\textbf{Design.} Five contrasts, each computed inside strata and averaged over those with both populations
(weights $n_Pn_Q/(n_P+n_Q)$, $n$ the flights of each side): $\G-\Sl$ pooled; $\G-\Sl$ and
$\A-\Sl$ (\emph{selection}) at the same site-day and hour; $\G-\A$ within the flight and at the same
site-day and hour (the \emph{group effect}, estimated in two ways that fail in opposite ways: the
first compares a flight with itself at another time of day, the second compares pilots at the same
time). Where the last two agree, the effect belongs to the group.

\begin{figure}[H]
\centering
\includegraphics[width=\textwidth]{decompose.pdf}
\caption{Each statistic (panel) for the five contrasts (rows), with 5--95\% bands. Shares of time in
points; lower height lost per km is a better glide.}
\label{fig:decompose}
\end{figure}

\begin{table}[H]
\centering\footnotesize\setlength{\tabcolsep}{3.5pt}
\begin{tabular}{@{}l c c c c c@{}}
\toprule
$\Delta$ & \makecell{$\G-\Sl$\\pooled} & \makecell{$\G-\Sl$\\same s-d-h} & \makecell{$\A-\Sl$\\same s-d-h} & \makecell{$\G-\A$\\same flight} & \makecell{$\G-\A$\\same s-d-h} \\
\midrule
$H_\mathrm{loc}(100~\mathrm{s})$ & \makecell{$+$0.008\\{\scriptsize[$+$0.007, $+$0.009]}} & \makecell{$+$0.002\\{\scriptsize[$+$0.001, $+$0.003]}} & \makecell{$+$0.005\\{\scriptsize[$+$0.004, $+$0.005]}} & \makecell{$-$0.009\\{\scriptsize[$-$0.009, $-$0.008]}} & \makecell{$+$0.002\\{\scriptsize[$+$0.001, $+$0.002]}} \\[2pt]
$H_\mathrm{loc}(1000~\mathrm{s})$ & \makecell{$+$0.055\\{\scriptsize[$+$0.053, $+$0.057]}} & \makecell{$+$0.041\\{\scriptsize[$+$0.038, $+$0.044]}} & \makecell{$+$0.003\\{\scriptsize[$+$0.001, $+$0.006]}} & \makecell{$+$0.037\\{\scriptsize[$+$0.035, $+$0.039]}} & \makecell{$+$0.043\\{\scriptsize[$+$0.040, $+$0.046]}} \\[2pt]
time in climb (\%) & \makecell{$+$2.20\\{\scriptsize[$+$1.96, $+$2.45]}} & \makecell{$+$3.43\\{\scriptsize[$+$3.19, $+$3.63]}} & \makecell{$+$0.43\\{\scriptsize[$+$0.25, $+$0.59]}} & \makecell{$+$5.71\\{\scriptsize[$+$5.52, $+$5.92]}} & \makecell{$+$3.28\\{\scriptsize[$+$3.11, $+$3.47]}} \\[2pt]
searches to climb (\%) & \makecell{$+$6.18\\{\scriptsize[$+$5.82, $+$6.46]}} & \makecell{$+$4.12\\{\scriptsize[$+$3.68, $+$4.54]}} & \makecell{$+$1.60\\{\scriptsize[$+$1.23, $+$1.97]}} & \makecell{$+$8.75\\{\scriptsize[$+$8.37, $+$9.15]}} & \makecell{$+$2.59\\{\scriptsize[$+$2.25, $+$2.99]}} \\[2pt]
climb rate (m/s) & \makecell{$+$0.040\\{\scriptsize[$+$0.030, $+$0.051]}} & \makecell{$+$0.031\\{\scriptsize[$+$0.025, $+$0.039]}} & \makecell{$+$0.007\\{\scriptsize[$+$0.001, $+$0.012]}} & \makecell{$+$0.024\\{\scriptsize[$+$0.019, $+$0.031]}} & \makecell{$+$0.029\\{\scriptsize[$+$0.025, $+$0.034]}} \\[2pt]
glide speed (km/h) & \makecell{$+$3.18\\{\scriptsize[$+$2.97, $+$3.39]}} & \makecell{$+$1.97\\{\scriptsize[$+$1.82, $+$2.10]}} & \makecell{$+$1.15\\{\scriptsize[$+$1.05, $+$1.27]}} & \makecell{$+$1.36\\{\scriptsize[$+$1.27, $+$1.46]}} & \makecell{$+$0.87\\{\scriptsize[$+$0.79, $+$0.98]}} \\[2pt]
height lost per km (m/km) & \makecell{$-$18.6\\{\scriptsize[$-$19.3, $-$17.8]}} & \makecell{$-$10.0\\{\scriptsize[$-$10.7, $-$9.1]}} & \makecell{$-$0.5\\{\scriptsize[$-$1.1, $+$0.3]}} & \makecell{$-$21.1\\{\scriptsize[$-$22.1, $-$20.1]}} & \makecell{$-$9.5\\{\scriptsize[$-$10.5, $-$8.6]}} \\[2pt]
cross-country speed (km/h) & \makecell{$+$2.81\\{\scriptsize[$+$2.59, $+$3.01]}} & \makecell{$+$2.25\\{\scriptsize[$+$2.09, $+$2.39]}} & \makecell{$+$1.31\\{\scriptsize[$+$1.21, $+$1.43]}} & \makecell{$+$0.70\\{\scriptsize[$+$0.60, $+$0.78]}} & \makecell{$+$1.29\\{\scriptsize[$+$1.16, $+$1.41]}} \\[2pt]
\bottomrule
\end{tabular}

\caption{The contrasts of Figure~\ref{fig:decompose}: difference and 5--95\% band (s-d-h: same site-day and hour).}
\label{tab:decompose}
\end{table}

\begin{reading}
\textbf{Reading.}
\begin{itemize}
  \item \textbf{Straighter progress over 1000~s} survives every control: $\Delta H_\mathrm{loc}(1000)$ is
        \DHlocThousandRaw{} pooled, \DHlocThousandFlight{} within the flight, \DHlocThousandGroupSdh{} at the same
        site-day and hour; selection only \DHlocThousandSel.
  \item \textbf{Nothing within one phase (100~s):} \DHlocHundredFlight{} within the flight,
        \DHlocHundredGroupSdh{} at the same place and hour; the pooled \DHlocHundredRaw{} was selection and composition.
  \item \textbf{More time climbing} (\DClimbFlight{}, \DClimbGroupSdh{} points) and \textbf{more searches that
        end in a climb} (\DFoundFlight{}, \DFoundGroupSdh{} points).
  \item \textbf{Climb rate is higher in a group} (\DRateFlight, \DRateGroupSdh{} m/s), not lower as in
        the pooled table: a composition effect (pooled ratios weight flights by their time in each population).
  \item \textbf{Faster and flatter glides:} \DSpeedFlight, \DSpeedGroupSdh{} km/h and \DSinkFlight, \DSinkGroupSdh{} m/km
        less height lost per km; selection adds \DSpeedSel{} km/h of glide speed (group pilots glide faster even alone).
  \item \textbf{Cross-country speed} is higher in a group once stratified: \DXcFlight, \DXcGroupSdh{} km/h.
\end{itemize}
\end{reading}

%--------------------------------------------------------------------------------------
\subsection{Size of the configuration}\label{sec:size}

The group time is cut where the configuration changes: every stretch has a fixed number of pilots $n$
(2, 3, 4, 5--6, 7+); a flight can contribute to several sizes. Two contrasts by size: each size minus the
same flights alone (within the flight), and each size minus configurations of 2 pilots (same site-day
and hour: does the size matter beyond a pair?).

\begin{figure}[H]
\centering
\includegraphics[width=0.56\textwidth]{size_hloc.pdf}
\caption{$H_\mathrm{loc}(\tau)$ over every lag by the size of the configuration (orange, darker with more pilots), S (blue) and, dashed, the same
flights alone for each size.}
\label{fig:size-hloc}
\end{figure}

\begin{table}[H]
\centering\footnotesize\setlength{\tabcolsep}{3pt}
\begin{tabular}{@{}l r c c c c c c c c c@{}}
\toprule
 & & \multicolumn{2}{c}{$H$} & \multicolumn{2}{c}{$H_\mathrm{loc}(1000~\mathrm{s})$} & & & & & \\
\cmidrule(lr){3-4}\cmidrule(lr){5-6}
pilots & flights & $\G_n$ & alone & $\G_n$ & alone & \makecell{time in\\climb (\%)} & \makecell{searches\\to climb (\%)} & \makecell{climb\\rate (m/s)} & \makecell{glide speed\\(km/h)} & \makecell{glide\\ratio} \\
\midrule
solo & 119\,129 & 0.899 &  & 0.837 &  & 30.8 & 43.3 & 1.22 & 30.1 & 10.4 \\
2 & 26\,411 & 0.917 & 0.914 & 0.889 & 0.861 & 32.7 & 49.0 & 1.27 & 33.2 & 12.5 \\
3 & 9\,480 & 0.923 & 0.919 & 0.906 & 0.867 & 33.6 & 49.9 & 1.26 & 33.5 & 13.1 \\
4 & 4\,073 & 0.924 & 0.920 & 0.910 & 0.865 & 33.6 & 51.0 & 1.25 & 33.5 & 13.7 \\
5--6 & 2\,686 & 0.925 & 0.922 & 0.915 & 0.870 & 33.1 & 50.1 & 1.23 & 33.1 & 14.1 \\
7+ & 1\,142 & 0.933 & 0.924 & 0.944 & 0.871 & 34.6 & 51.7 & 1.20 & 34.6 & 15.9 \\
\bottomrule
\end{tabular}

\caption{Paragliders, pooled by the size of the configuration (``alone'': the same flights alone).}
\label{tab:size}
\end{table}

\begin{figure}[H]
\centering
\includegraphics[width=\textwidth]{decompose_size.pdf}
\caption{Orange: $\G_n-\A$ within the flight. Blue: $\G_n-\G_2$ at the same site-day and hour.}
\label{fig:size}
\end{figure}

\begin{reading}
\textbf{Reading.} Pooled, $H$ and $H_\mathrm{loc}$ rise with $n$ while the same flights alone stay flat (Fig.~\ref{fig:size-hloc}, Table~\ref{tab:size}). Within the flight the straight-progress effect grows with $n$ (\DHlocThousandFlightTwo{} with 2 pilots,
\DHlocThousandFlightSeven{} with 7+), but at the same site-day and hour configurations of 3 to 6 pilots do not differ from pairs
($\Delta H_\mathrm{loc}(1000)$ between \DHlocThousandMiddleMin{} and \DHlocThousandMiddleMax, bands across zero): the rise
of the pooled curves belongs to the days and sites where large groups form. Only 7+ stands apart
(\DHlocThousandSevenTwoSdh{} \DHlocThousandSevenTwoSdhBand, cross-country speed \DXcSevenTwoSdh{} km/h), a rare
configuration (3\% of group time) of unknown nature. Time in climb and the flatter glide do grow with $n$ at the same place
and hour (\DClimbSizeMin{} to \DClimbSizeMax{} points; \DSinkSizeMin{} to \DSinkSizeMax{} m/km).
\end{reading}

%--------------------------------------------------------------------------------------
\subsection{Leader and follower}\label{sec:mechanisms}

In configurations of two pilots the signed median delay of the pair interval says who passes first: at each instant,
the delay at which the other pilot passes closest, and its median over the instants with neither pilot in a search
(leader: the other
passes at least 15~s later; \Abreast\% abreast, median delay \DelayMedian~s). The two are compared within the
configuration (Appendix~\ref{app:leaders}). How groups form and break, whether the same pilots lead, and when groups
fly are in section 7 of the note \emph{Group flights in free flight}.

\begin{reading}
\textbf{Reading.}
The pilot ahead flies a little faster (glide speed \LeadSpeed{} km/h, climb rate \LeadRate{} m/s) but does not find thermals more
often (searches ending in a climb \LeadFound{} points, height lost per km \LeadSink{} m/km): within a pair the follower
shows no sign of using the leader to find the next thermal. Group time is concentrated around midday and its time alone
falls earlier and later (section 7 of the note): this is why $\G-\A$ within the flight must be read with the
same-site-day-hour estimate.
\end{reading}
